\section{Introduction}
\label{sec:intro}

\subsection{The conjecture}
\label{ssec:conjecture}

A finite family $\F$ of finite sets is \emph{union-closed} if
$A,B\in\F$ implies $A\cup B\in\F$. The Union-Closed Sets Conjecture,
attributed to Péter Frankl around 1979, asserts:

\begin{conjecture}[Frankl, union-closed form]
\label{conj:frankl}
For every finite union-closed family $\F$ with $\F\neq\{\emptyset\}$,
there is an element $x$ of the ground set belonging to at least half
the members of $\F$.
\end{conjecture}

The conjecture is elementary to state and notoriously resistant. After
four decades in which only sporadic special cases were known --- see
the survey of Bruhn and Schaudt~\cite{bruhn-schaudt} --- the first
unconditional constant lower bound was obtained by
Gilmer~\cite{gilmer} in 2022 via an entropy argument, showing that some
element lies in at least a $0.01$ fraction of the family; the entropy
method was then sharpened by Sawin and others to a fraction of
$(3-\sqrt5)/2\approx0.382$~\cite{sawin}, which is also where the
entropy method, in its decoupled form, provably stalls. The threshold
$\tfrac12$ remains open.

This paper consolidates a research program --- recorded across fourteen
internal notes, here referred to as the \emph{UC1--UC14 ledger} --- that
attacks Conjecture~\ref{conj:frankl} from a structural rather than an
analytic direction. The program has two arcs. The first
(Sections~\ref{sec:setup}--\ref{sec:combinatorial}) is a native
combinatorial reduction: it rewrites Frankl's conjecture as an
increasingly sharp sequence of explicit finitary inequalities. The
second (Sections~\ref{sec:cohomology}--\ref{sec:obstruction}) is a
cohomological program: it builds an equivariant cohomology theory of
intersection-closed families and aims to close Frankl by a two-part
contradiction on a distinguished obstruction class.

\subsection{What this paper is, and what it is honest about}
\label{ssec:honesty}

\textbf{The Union-Closed Sets Conjecture is not proved in this paper,
and the program does not currently prove it.} The UC1--UC14 ledger is a
chain of internally generated, \emph{unrefereed} research notes; at
various points it carried internal verdict banners claiming a complete
proof. An audit pass over the ledger (recorded in four further notes,
consolidated here in Section~\ref{sec:proof-state}) found that those
banners overstated the state. The present document is the corrected
consolidation: it presents the genuine mathematical content of the
program --- which is substantial --- and marks, explicitly and at each
load-bearing step, what is proved, what is reduced, what is
axiomatised, and what is false.

To that end we use three referee-facing environments throughout. A
\emph{Status} paragraph records the verified standing of the
immediately preceding claim. An \emph{Open residual} environment states
a precisely isolated open problem. A \emph{Correction} environment
records a place where the earlier ledger was wrong and explains why, so
that the error is not re-propagated. The reader who wants only the
bottom line should read Section~\ref{ssec:status-summary} below and
then Section~\ref{sec:proof-state}.

A companion machine-checked formalisation exists (in the Lean~4 proof
assistant, against \texttt{mathlib}); its current state is reported
honestly in Section~\ref{ssec:lean-state}. This paper follows a
\emph{math-before-Lean} discipline: the careful mathematical account is
primary, and the formalisation is treated as a witness to be reported
on, not as a source of authority. Where the formalisation and the paper
disagree, Section~\ref{sec:proof-state} says so.

\subsection{The structural analogy}
\label{ssec:analogy}

The program's organising idea is that Frankl's conjecture is an
instance of a recurring pattern --- \emph{compatibility constraints
plus local symmetry force a global bias} --- shared with the
$1/3$--$2/3$ (Kahn--Saks) conjecture for finite posets. The two
problems are not formally reducible to one another, but they admit a
common structural template:

\begin{center}
\begin{tabular}{ll}
\toprule
$1/3$--$2/3$ (Kahn--Saks) & Union-Closed (Frankl)\\
\midrule
linear extensions $L(P)$ & members $A\in\F$\\
adjacent transpositions & coordinate transports / toggles\\
Bruhat / weak-order geometry & cubical geometry of the Boolean cube\\
order constraints & union closure\\
pair-balance & coordinate-frequency balance\\
conductance obstruction & pairing obstruction\\
\bottomrule
\end{tabular}
\end{center}

In both settings the ambient geometry is highly symmetric, the
admissible moves preserve only part of that symmetry, local
cancellation looks possible, and global cancellation is conjecturally
obstructed. The program's wager is that the obstruction is
cohomological: it arises because a family of local symmetries fails to
globalise coherently. Section~\ref{sec:setup} makes the dictionary
precise; the manifesto note (UC1) records that, while the framework
transfers cleanly, \emph{neither} of the two proof engines used on the
$1/3$--$2/3$ side transfers, so the program had to build a native
engine. Sections~\ref{sec:combinatorial}
and~\ref{sec:cohomology}--\ref{sec:obstruction} are the two native
engines it built.

\subsection{Status summary}
\label{ssec:status-summary}

We summarise the standing of every load-bearing component. The
detailed account is Section~\ref{sec:proof-state}; the cross-reference
ledger is Appendix~\ref{sec:appendix-status}.

\begin{description}[leftmargin=1.4em,style=nextline]

\item[The combinatorial arc (UC1--UC8), Section~\ref{sec:combinatorial}.]
\emph{Genuine partial results; a complete reduction; not a proof.} The
arc rewrites Frankl as a single explicit integer inequality on the
trace data of a minimal generator (Theorems~\ref{thm:trace-partition}
and~\ref{thm:boundary-conservation}) and settles that
inequality on a strictly growing hierarchy of structured
sub-families. The cases of generator size $\le2$ --- equivalently
families with a set of size $\le2$ --- are settled
\emph{unconditionally and natively} (Theorem~\ref{thm:trace-partition}),
recovering classically known results. No structural class settled so
far exhausts union-closure, and the program proves a genuine negative
result --- no per-fibre sign theorem can exist
(Proposition~\ref{prop:no-per-fibre}) --- showing the residual is
irreducibly a global statement.

\item[The cohomological arc, trace-injectivity, Section~\ref{sec:cohomology}.]
\emph{One genuinely unconditional theorem.} The doubling functor and
the matched-cylinder null-homotopy yield an honest, unconditional
proof that the canonical trace map vanishes on the top Walsh isotype
(Theorem~\ref{thm:trace-injectivity}). This is the UC12 result; it does
not depend on any of the disputed steps below.

\item[$Y_2$ --- the obstruction-class landing, Section~\ref{ssec:y2}.]
\emph{Mechanism sound.} The claim that the obstruction class lands in
the level-$1$ Walsh isotypes rests on a representation-theoretic
mechanism --- an additive equivariant map preserves Walsh isotype, and
only a cup product can shift isotype level --- which has been
machine-checked (Theorem~\ref{thm:y2-mechanism}). The competing earlier
claim (a top-isotype landing) is machine-checked to be impossible. A
caveat remains on the faithfulness of one input; see
Section~\ref{ssec:y2}.

\item[$Y_1$ --- the level-$1$ Walsh vanishing, Section~\ref{ssec:y1}.]
\emph{Open. Not proved.} The previously circulated proof
(UC13 \S4.5) is logically invalid: it conflates exactness of a trace
map with vanishing of cohomology (Correction~\ref{corr:trace-exact}).
The honest re-derivation by cofiber long-exact-sequence and induction
\emph{validly reduces} $Y_1$ to a single named open computation, but
does not close it (Open residual~\ref{res:y1}). $Y_1$ is
true-but-unproven.

\item[UC14 ``$R2$''. ]
\emph{False.} The chain-isomorphism lemma advertised as repairing the
$Y_1$ gap is provably false; it transports the program's own nonzero
sphere class into a group it requires to vanish
(Correction~\ref{corr:r2}).

\item[The Lean headline, Section~\ref{ssec:lean-state}.]
\emph{Axiom-dependent.} The formalised theorem \texttt{Frankl\_Holds}
rests on one named project axiom, \texttt{case3\_ss\_obstruction\_%
paper\_axiom}, which --- combined with a proved companion theorem ---
is propositionally equivalent to Frankl itself
(Section~\ref{ssec:axiom-is-frankl}). The genuinely unconditional
machine-checked content is the concrete full-powerset families at
$n=3$ and $n=4$ (Section~\ref{ssec:lean-unconditional}).

\end{description}

\noindent The single sentence that captures the state: \emph{the
program reduces Frankl, on the cohomological side, to one named
homotopy-colimit computation, and on the combinatorial side to one
named integer inequality; both reductions are correct; neither
residual is closed.}

\subsection{Organisation}
\label{ssec:organisation}

Section~\ref{sec:setup} fixes notation, proves the transport identity,
gives the four equivalent forms of the obstruction, and introduces the
intersection-closed dialect and the Walsh picture.
Section~\ref{sec:combinatorial} is the combinatorial reduction arc
(UC1--UC8). Section~\ref{sec:cohomology} builds the equivariant
cohomology of intersection-closed families (UC9--UC12) and proves the
trace-injectivity theorem. Section~\ref{sec:obstruction} assembles the
\v{C}ech obstruction class and states the intended two-part
contradiction (UC11, UC13). Section~\ref{sec:proof-state} is the honest
audit: it states and proves the corrections, isolates $Y_1$ as open,
and reports the Lean state. Section~\ref{sec:residual} states the named
open residuals and the outlook. Appendix~\ref{sec:appendix-notation}
is a notation index; Appendix~\ref{sec:appendix-status} is the
component-by-component status ledger.

\subsection{Provenance and scope}
\label{ssec:provenance}

The mathematical content consolidated here was developed in the notes
UC1--UC14 and audited in four further notes. This paper is a
\emph{consolidation and correction}, not new research: it does not
attempt to close either named residual. Its purpose is to render the
program's genuine content in one referee-readable document and to
record its honest state so that the known defects are not
re-propagated. Two structural inputs are imported from a sibling
program on the $1/3$--$2/3$ conjecture and used as black boxes; they
are flagged at the point of use (Section~\ref{ssec:u1-wall} and
Remark~\ref{rem:external-input}). A referee assessing the cohomological
arc should treat those imports as hypotheses to be checked against
their stated sources.
